% =============================================================================
% Pflichtenheft.tex — Bachelor Thesis Template (ZHAW / IMS)
% =============================================================================
\documentclass[a4paper,11pt,bibliography=totoc,listof=totoc]{scrreprt}

\usepackage[english]{babel}

% bastyl loads glossaries-extra internally -> pass required option BEFORE bastyl
\PassOptionsToPackage{automake=immediate}{glossaries-extra}
\usepackage{bastyl}

% bastyl/glossaries setup expects this
\makeglossaries

% =============================================================================
% Thesis Metadata
% =============================================================================
\renewcommand{\thesistitle}{Betaflight -- Offline Controller Tuning}
\renewcommand{\thesissubtitle}{%
  Data-Driven Tuning of Rate and Altitude Controllers for FPV Drones%
}
\renewcommand{\thesisauthors}{Yuri Bianchi, Dario Jurietti and Janick Dort}
\renewcommand{\thesisdate}{\today}
\renewcommand{\thesisheader}{Betaflight -- Offline Controller Tuning}
\renewcommand{\thesisfooter}{ZHAW / IMS}

\hypersetup{
  pdftitle    = {Betaflight -- Offline Controller Tuning Pflichtenheft},
  pdfauthor   = {Yuri Bianchi, Dario Jurietti, Janick Dort},
  pdfsubject  = {Bachelor Thesis, ZHAW School of Engineering},
  pdfkeywords = {Betaflight, PID tuning, system identification, FPV, chirp},
}

\begin{document}

\section{Introduction}

First-Person-View (FPV) racing drones require precise and responsive control 
in order to achieve stable flight behaviour, suppress disturbances such as 
propeller wash and maintain agility during rapid manoeuvres. The quality of this 
behaviour strongly depends on the tuning of the flight controller’s feedback loops. 
In particular, the parameters of the angular rate controller determine how quickly 
and accurately the drone reacts to pilot inputs and external disturbances.

In practice, controller tuning is still predominantly performed manually through 
iterative trial-and-error procedures. This process requires extensive flight 
testing and considerable piloting experience, making it both time-consuming and 
difficult to reproduce. Even though modern flight controller firmware such as 
Betaflight offers extensive configuration possibilities and detailed flight data 
logging, efficient and systematic tuning remains a significant challenge.

A previous project introduced a proof-of-concept approach for offline tuning of 
angular rate controllers using recorded flight data and system identification 
techniques. Since then, this concept has been further developed beyond the initial 
proof of concept and forms the basis for the present work.

The objective of this bachelor thesis is to further refine and extend this approach 
into a modular offline tuning framework for multiple control loops used in 
Betaflight-based flight controllers. In addition to the already developed angular 
rate controller tuner, new modules will be introduced for tuning the angle 
controller, the altitude hold controller, and the position controller. These 
additions aim to extend the tuning framework beyond the inner rate control loop 
and provide tools for analysing and tuning higher-level control layers.

The framework will be implemented primarily in MATLAB and Python and will provide 
tools for analysing flight logs, estimating system dynamics and generating 
controller parameters without requiring repeated flight tests. The software will 
be developed as an open-source library and accompanied by comprehensive 
documentation to ensure accessibility and usability for the FPV community.

By providing a structured, data-driven approach to controller tuning, this 
project aims to simplify the tuning process, improve flight performance and 
reduce the reliance on manual trial-and-error methods.

\section{Goal / Task Description}

The goal of this bachlelor thesis is to develop a modular offline tuning framework
for multiple control loops used in Betaflight-based flight controllers. Newly,
the angle controller and the altitude hold controller should be added to the
already developed angular rate controller tuner. If time allows, a position 
controller tuner should be also implemented as well. The framework should be 
implemented primarily in MATLAB and Python. 

To achieve this goal, the firmware of the flight controller will be modified 
to add the chirp signal to the input signal of the angle controller, the 
altitude hold controller and the position controller.

\section{Specification}
\subsection{Technical Specification}

For the implementation of the offline tuning framework, the following technical
specifications are defined:
\begin{itemize}
    \item The framework should be implemented in MATLAB and Python.
    \item The framework should be modular, allowing for the addition of new tuning modules for different control loops.
    \item The framework should provide tools for analysing flight logs, estimating system dynamics and generating controller parameters.
    \item The framework should be developed as an open-source library and accompanied by comprehensive documentation.
\end{itemize}

\subsection{Functional Specification}
The functional specifications for the offline tuning framework are as follows:
\begin{itemize}
    \item The framework should be able to process flight logs recorded with Betaflight firmware.
    \item The framework should be able to estimate the system dynamics of the angle controller, 
    the altitude hold controller and the position controller based on the recorded flight data.
    \item The framework should be able to generate controller parameters for the angle controller, 
    the altitude hold controller and the position controller based on the estimated system dynamics.
    \item The framework should provide a user-friendly interface for configuring and running the tuning process.
\end{itemize}

\subsection{Boundary Conditions}
The following boundary conditions are defined for the offline tuning framework:
\begin{itemize}
    \item The framework should be compatible with Betaflight firmware version 
    \item The framework should be designed for use with FPV racing drones equipped with Betaflight-based flight controllers.
\end{itemize}

\section{Costs}
The costs associated with the development of the offline tuning framework are estimated as follows:
\begin{itemize}
    \item Additional GPS modules for testing the position controller tuner
\end{itemize}

\section{Work Packages}
The project is divided into several work packages that structure the development 
and validation of the offline tuning framework. 

\end{document}